
This paper set out to measure whether ML/NLP and HCI alignment researchers cite each other symmetrically within the huashen218 bidirectional alignment corpus. The measurement infrastructure was found to be absent and was constructed. In building it, a methodological discovery was made: the default venue-string classification approach used in most bibliometric tools achieves only 12\% label coverage on interdisciplinary alignment corpora indexed by S2AG; FoS-primary classification achieves 100\%.

The ''corpus'' accessible through the public GitHub reading list contains 49 extractable paper identifiers, not the approximately 400 papers described in the full systematic review. In the 33 papers resolved via S2AG, the directional asymmetry consistent with H-CitAsym was already present: a directed citation ratio of 0.107, with HCI alignment papers allocating 100\% of within-corpus outgoing citations to ML/NLP work. Both pre-registered gate thresholds were not met, and no statistical tests were executed. All quantitative predictions are INCONCLUSIVE.

Four contributions are reported: (1) a validated S2AG bibliometric pipeline (23/23 unit tests, fully cached); (2) FoS-primary classification achieving 100\% versus 12\% label coverage on interdisciplinary alignment corpora; (3) the first directional citation measurement in the huashen218 corpus (ratio $\approx$ 0.107, directionally consistent with H-CitAsym but not statistically confirmable at N = 9 cross-group edges); and (4) a documented methodological finding that curated reading lists yield approximately 12\% of the described corpus size --- programmatic S2AG citation-of-citation search is the viable reconstruction approach.

The immediate next step is corpus augmentation: using S2AG's endpoint for papers citing Shen et al. [2024] to recover 200+ papers and enable full statistical testing. All pipeline components are implemented and validated; corpus augmentation is the remaining step before statistical testing of P1, P2, and P3 becomes feasible.

